\chapter{Introduction}
\label{chap:introduction}

Monte Carlo simulation is a standard technique for pricing derivatives whose payoffs or dynamics are inconvenient to evaluate analytically. In quantitative finance it is rarely enough to compute only a price: desks also require sensitivities, or Greeks, for hedging, risk reporting and scenario analysis. Automatic differentiation (AD) offers a systematic way to obtain these sensitivities from simulation code, but it changes the performance profile of the computation. Runtime depends not only on the number of paths, but also on the payoff, time discretisation, implementation language, compiler behaviour, AD mode, threading, numerical libraries and the hardware on which the benchmark is run.

This creates a practical benchmarking problem. A quantitative developer choosing between a Python/NumPy implementation, JAX/XLA, a C++ extension or a Rust extension cannot rely on language reputation alone. A fast engine for a simple European option may not be the best choice for a path-dependent or local-volatility workload, and an engine that is attractive on a laptop may have a different cost-performance profile on a cloud instance. Exhaustive benchmarking can resolve this uncertainty, but it becomes expensive when the configuration grid includes multiple workloads, engines, path counts, AD modes and cloud machines.

\section{Problem and Gap}

The project addresses the gap between small, ad hoc microbenchmarks and a reproducible benchmarking process suitable for Monte Carlo workloads with AD. Existing financial Monte Carlo references establish the pricing methodology, and AD literature explains the derivative computation model, but neither directly answers the engineering question considered here: how should one compare MC+AD implementations across engines and cloud instances when full-grid benchmarking is costly?

The goal was therefore not to build a production pricing library. Instead, the project set out to build and evaluate a reproducible benchmarking framework that can:
\begin{itemize}
    \item express comparable option-pricing workloads across several engines;
    \item validate numerical results where analytical or cross-engine checks are available;
    \item persist timings, metadata and cloud cost information in a queryable form;
    \item compare engine performance, AD overhead and cloud cost-performance; and
    \item reduce the number of expensive full-scale benchmark evaluations using cheap probes.
\end{itemize}

\section{Research Questions}

The report is organised around five research questions:
\begin{enumerate}
    \item \textbf{RQ1: Correctness and reproducibility.} Can the framework run repeatable MC+AD experiments while retaining enough metadata to audit the results?
    \item \textbf{RQ2: Engine performance.} How do NumPy, JAX/XLA, C++/OpenMP and Rust/Rayon implementations compare across the evaluated workloads?
    \item \textbf{RQ3: AD overhead.} What runtime overhead is introduced by JAX forward- and reverse-mode AD for the supported workloads?
    \item \textbf{RQ4: Cloud cost-performance.} How does runtime translate into per-run cost on priced cloud instances?
    \item \textbf{RQ5: Budget-aware profiling.} Can a Successive Halving style profiler reduce full benchmark evaluations while preserving the choices made by a full grid search?
\end{enumerate}

\section{Contributions}

The main contributions are:
\begin{itemize}
    \item a modular benchmarking framework with typed workload configurations, engine abstraction, repeated timing, deterministic seeding, SQLite storage and CSV export;
    \item implementations of European, Asian and parametric local-volatility option workloads across NumPy, JAX, C++ and Rust engines where supported;
    \item integration of cloud metadata and cost estimates into benchmark results;
    \item an old probe-based profiler and a Successive Halving Algorithm (SHA) profiler for reducing full-scale benchmark runs; and
    \item an empirical evaluation using real result files, including correctness checks, engine comparisons, AD overheads, cloud cost-performance, Pareto analysis and profiler selection quality.
\end{itemize}

The strongest result is the final cloud profiling experiment, \code{sha_cloud_profiler_v4}. On an \code{n2-standard-4} instance in \code{europe-west1-b}, the complete grid contained 16 configurations and 48 full benchmark runs. SHA selected 24 full runs, saving 50.0\% of the full-grid runs while recovering 100.0\% of the full-grid Pareto frontier and incurring 0.0\% runtime and cost regret relative to the grid best. The best runtime configuration, \code{rust/european/none}, took 0.54 ms at \(M=100{,}000\) and was selected by the full grid, the old profiler and SHA.

\section{Scope}

The evaluation is intentionally scoped. It covers CPU-focused local and Google Cloud experiments, with JAX used for AD and C++/Rust used for no-AD compiled engines. GPU execution, Mojo, distributed execution, hardware performance counters, wider cloud regions and a systematic BLAS comparison were not completed and are treated as future work. This scope is important: the results suggest useful engineering choices for the measured experiments, but they should not be read as universal performance laws.

\section{Report Structure}

Chapter~\ref{chap:background} summarises the relevant background in Monte Carlo pricing, AD, benchmarking and Successive Halving. Chapter~\ref{chap:design} describes the framework design. Chapter~\ref{chap:implementation} explains the implemented system and its storage, cloud and profiling components. Chapter~\ref{chap:experiments} defines the experiments and metrics. Chapter~\ref{chap:evaluation} presents and discusses the results. Chapter~\ref{chap:conclusion} concludes with limitations and future work.
